The model was developed in an existing Agent-Based Modelling Software Python library called \textit{MESA} \citep{Mesa2025}, a methodological choice strongly supported by recent reviews of computational tools in solid waste management \citep{TianReview2024, Ma2023}. Its architecture consisted of one single ABM environment that contained 7 distinct \texttt{BarangayAgent} objects, which in turn contained their respective populations of \texttt{HouseholdAgent} objects.

The simulation utilized a single, discrete ABM environment to contain its operational entities and manage its temporal structure. Time progression was segmented into recurring quarterly loops, establishing the fundamental unit of analysis and adaptation for the governance agents. The architecture was strictly hierarchical: it contained seven distinct \texttt{BarangayAgent} objects in the upper layer, which encapsulated their respective populations of \texttt{HouseholdAgent} objects in the lower layer. This multi-level approach was essential for accurately capturing the flow of mandates and the heterogeneity between different governance units in socio-environmental systems \citep{Brugiere2022}.

\begin{figure}[h]
    \centering
    % Ensure you have uploaded the file 'multi-agent.png' to your project
    \includegraphics[width=0.8\linewidth]{multi-agent.png}
    \caption{Multi-Level Agent-Based Modeling Diagram}
    \label{fig:multi_level_abm}
\end{figure}

% Define a custom bullet for tables to save space and look cleaner
\newcommand{\tabitem}{~~\llap{\textbullet}~~}

\begin{table}[htbp]
    \centering
    \small
    \caption{Agent Role Mapping and Architecture}
    \label{tab:agent_roles_v2}
    \renewcommand{\arraystretch}{1.5} 
    \begin{tabularx}{\textwidth}{ 
        >{\RaggedRight\hsize=0.15\hsize}X 
        >{\RaggedRight\hsize=0.12\hsize}X 
        >{\RaggedRight\hsize=0.29\hsize}X 
        >{\RaggedRight\hsize=0.22\hsize}X 
        >{\RaggedRight\hsize=0.22\hsize}X 
    }
        \toprule
        \textbf{Agent} & \textbf{Type} & \textbf{Input / Observation} & \textbf{Action Space} & \textbf{Objective} \\
        \midrule
        \textbf{Household} & ABM (Env.) & 
        \tabitem Policy: Incentive \& Fine amounts \newline
        \tabitem Social: Neighbor compliance \newline & 
        Decide: Segregate or Do Not Segregate (Binary) & 
        Maximize Utility: Minimize ``Cost of Effort'' vs. Fear of Penalty. \\
        \textbf{Barangay} & ABM (Impl.) & 
        \tabitem Budget: Quarterly funds \newline
        \tabitem Population: Count \& income profile & 
        Execute: Deploy Enforcers or distribute rewards & 
        Report State: Aggregate data and transmit to LGU. \\
        \textbf{LGU (Municipal)} & DRL / AI (Brain) & 
        \tabitem State: Compliance of 7 Barangays \newline
        \tabitem Fiscal: Remaining Budget \newline
        \tabitem Political: Political Capital & 
        Allocate: Distribute budget to (1) IEC, (2) Incentives, (3) Enforcement. & 
        Maximize Reward: Balance compliance with sustainability. \\
        \textbf{Enforcement} & ABM (Ops.) & 
        \tabitem Target: Assigned Barangay \newline
        \tabitem Constraint: Working hours & 
        Inspect: Randomly check households & 
        Detect: Trigger ``Perceived Probability of Detection''. \\
        \bottomrule
    \end{tabularx}
\end{table}

The \textit{BarangayAgents} embodied the Local Government Unit (LGU) implementation layer, responsible for fiscal management (fund allocation), policy formulation and implementation (e.g., establishing waste segregation mandates), and personnel management (deployment of enforcement staff). This aligned with the decentralized governance structure where local units were the primary implementers of environmental law \citep{Nishimura2022}. The \textit{BarangayAgents} were the key recipients of the feedback mechanism, utilizing the aggregated compliance rate as an observational State and a prescriptive Reward signal for its internal Deep Reinforcement Learning (DRL) mechanism \citep{Jimenez2025}.

The \textit{EnforcementAgents} functioned as the operational execution arm of the LGU, tasked with the stochastic monitoring of compliance. Constrained by the budget allocation determined by the \textit{BarangayAgents}, these agents conducted random inspections to verify waste segregation at the source, effectively operationalizing the "No Segregation, No Collection" policy \citep{Badua2022}. Their presence directly manipulated the "perceived intensity of enforcement" variable within the simulation, transforming abstract policy mandates into tangible risks for non-compliant households \citep{Carpio2025}. This agent layer introduced the critical dynamic of detection probability: while higher enforcement density significantly deterred non-compliance ("green criminology"), it incurred higher operational costs, forcing the system to balance fiscal sustainability with strict environmental policing \citep{Chen2023, Ibanez2022}.

The \textit{HouseholdAgents} represented the citizenry and were geographically and administratively segregated by their assigned barangay. Household behavior was modeled through a process of rational bounded observation, where agents assessed three primary exogenous factors: the clarity and strictness of LGU policies, the perceived intensity of local enforcement (driven by the \textit{EnforcementAgents}), and the compliance behavior of neighboring households (social norms) \citep{Meng2018, Liao2024}. This assessment culminated in an independent segregation decision, which dictated the agent's contribution to the collective compliance level. The quarterly feedback loop aggregated these decisions into seven discrete \textit{BarangayComplianceRate} values, which were then transmitted back to the LGU-DRL Agent, closing the adaptive system loop and facilitating the study of policy response to citizen behavior \citep{Ceschi2021, Fontaine2024}.

\begin{figure}[h]
    \centering
    % Ensure you have uploaded the file 'simulation_ui' to your project
    \includegraphics[width=0.8\linewidth]{simulation_ui.jpg}
    \caption{Bacolod Multi-View Simulation UI}
    \label{fig:simulatio_ui}
\end{figure}

\subsection{Formalization as a Markov Decision Process}

To facilitate the autonomous discovery of optimal strategies in the absence of longitudinal historical data, the interaction between the Local Government Unit (LGU) and the household agents was formalized as a finite-horizon Markov Decision Process (MDP). This mathematical framework transitioned the Agent-Based Model from a descriptive simulation into a prescriptive optimization environment, where the LGU-DRL agent iteratively interacted with a stochastic state space to maximize a cumulative utility function \citep{Kompella2020}. The MDP is formally defined by the quintuple $\mathcal{M} = \langle S, A, P, R, \gamma \rangle$, representing the state space, action space, transition probability, reward signal, and discount factor, respectively.

The State Space ($S$) is characterized as a multi-dimensional, continuous vector $S_t \in \mathbb{R}^{10}$, which provides the agent with a comprehensive snapshot of the municipal environment at each decision epoch. This vector encompasses the mean compliance rates across the seven constituent barangays, the normalized remaining budget for the fiscal year, and the simulation time index to account for seasonal or temporal variations in waste generation. By including these macroscopic variables, the agent can perceive the heterogeneous social landscape of the municipality, allowing for localized rather than monolithic policy responses \citep{HaMinh2025}. Correspondingly, the Action Space ($A$) is defined as a continuous vector $A_t \in \mathbb{R}^{21}$. This high-dimensional space represents the granular allocation of the quarterly budget across the three primary policy levers—Enforcement, Incentives, and Information Education Campaigns (IEC)—for each individual barangay. The choice of a continuous action space over a discrete one is critical for achieving "policy precision," as it enables the Reinforcement Learning agent to fine-tune resource intensity to the specific tipping points of household behavioral change \citep{Rajesh2025}.

The system dynamics are governed by the Transition Function ($P$), denoted as $P(S_{t+1} | S_t, A_t)$. Within this framework, the ABM serves as the implicit transition engine; unlike classical control problems where transitions are defined by explicit differential equations, the state evolution here emerges from the aggregate, non-linear behaviors of heterogeneous household agents updating their internal utility based on the Theory of Planned Behavior (TPB). Consequently, the simulation acts as a generative model where the subsequent state $S_{t+1}$ is a stochastic manifestation of the community’s collective response to the LGU's policy intervention $A_t$ \citep{TianReview2024, Jimenez2025}. To evaluate the efficacy of these interventions, the Reward Function ($R$) provides a scalar feedback signal $R_t$ that balances the LGU’s multi-objective mandates. This function is designed to reward high aggregate compliance while penalizing fiscal overspending and the social friction associated with excessive punitive enforcement \citep{Dey2025}. Finally, a Discount Factor of $\gamma = 0.99$ is applied to the objective function, ensuring that the DRL agent prioritizes strategies that foster long-term, sustained behavioral shifts rather than pursuing ephemeral, short-term compliance spikes that may deplete the annual budget prematurely.

\subsection{Household Agent Design}

The decision-making architecture of the \texttt{HouseholdAgent} is governed by a dynamic utility function grounded in the Theory of Planned Behavior (TPB), which posits that an individual's intention to perform a behavior is a function of their attitude, subjective norms, and perceived behavioral control \citep{Ceschi2021}. In contrast to static behavioral models, this framework incorporates time-variant weights ($w(t)$) for psychological constructs, allowing agent behavior to evolve non-linearly in response to LGU interventions and social feedback loops \citep{Taraghi2025, Ma2023}. The core decision logic is represented by the utility of segregation ($U_{\mathrm{segregate}}$):

\begin{equation}
    U_{\mathrm{segregate}} = (w_A(t) A) + (w_{SN}(t) SN_{\mathrm{local}}) + (w_{PBC}(t)  PBC) - C_{\mathrm{Net}} + \epsilon
    \label{eq:utility_function}
\end{equation}

In this formulation, $w_A(t)$ represents the temporal evolution of the agent's internal valuation of segregation, functioning as a decay model that increases in response to Information, Education, and Communication (IEC) investment and decays stochastically in the absence of reinforcement to simulate ``public forgetting'' \citep{Trushna2024}. The social component, $SN_{\mathrm{local}}$, is an endogenous variable derived from the observed compliance rate of the agent's immediate spatial neighborhood (radius $r$), capturing the effects of social pressure and observational learning \citep{Meng2018, Liao2024}. To account for the inherent uncertainty in human decision-making, the term $\epsilon$ introduces stochastic noise, ensuring the model reflects real-world behavioral variance \citep{Subedi2025}.

A critical innovation in this model is the calculation of $C_{\mathrm{Net}}$, the net behavioral cost, which characterizes the perceived friction of compliance. This variable is defined as:

\begin{equation}
    C_{\mathrm{Net}} = C_{\mathrm{Effort}} + (\gamma C_{\mathrm{Monetary}}) - (\gamma I) - (\gamma F P_{\mathrm{Detection}})
    \label{eq:net_disutility}
\end{equation}

\noindent This cost-benefit sub-model integrates both physical and financial barriers, defined as follows:

\begin{description}
    \item[$\bm{C_{\mathrm{Effort}}}$:] The physical hassle associated with the washing, sorting, and storage of waste.
    \item[$\bm{C_{\mathrm{Monetary}}}$:] Tangible financial expenses, such as the procurement of color-coded bins or sacks.
    \item[$\bm{\gamma}$ (Income Sensitivity):] A weighting factor derived from household income level; for lower-income households, $\gamma > 1$ as financial levers ($I$ and $F$) carry greater psychological weight, whereas $\gamma < 1$ for higher-income households.
    \item[$\bm{I}$ and $\bm{F}$:] The objective magnitudes of monetary incentives and punitive fines, respectively.
    \item[$\bm{P_{\mathrm{Detection}}}$:] The likelihood of enforcement detection, which is dynamically linked to the LGU's resource allocation for enforcement personnel.
\end{description}

Furthermore, the model captures complex psychological feedbacks such as ``psychological reactance.'' While LGU investments in IEC campaigns generally improve Attitude ($A$) and Subjective Norms ($SN$) \citep{Trushna2024}, the model stipulates that if enforcement intensity crosses a specific threshold, Attitude ($A$) may paradoxically decrease. This inverse reaction reflects the tendency of individuals to resist coercive mandates when they perceive a loss of autonomy, a concept central to Nudge Theory applications in waste management \citep{LoanBalanay2023}. Finally, Perceived Behavioral Control ($PBC$) is calibrated by the agent's internal psychological friction and self-efficacy regarding segregation. This ensures the simulation acknowledges that even agents with positive attitudes and strong social pressure may fail to comply if they perceive the behavior as too mentally taxing, overwhelmingly complex, or highly disruptive to their established daily routines.
\subsection{Behavioral Hysteresis and the Internalization of Social Norms}

While the baseline \texttt{HouseholdAgent} architecture operates on the Theory of Planned Behavior (TPB), standard models often exhibit unrealistic behavioral decay where compliance collapses immediately upon the cessation of external stimuli. To more accurately simulate the real-world stability of established habits, this study incorporates a \textit{Norm Internalization Mechanism}. This mechanism addresses the phenomenon of ``Cultural Inertia,'' where a behavior transitions from a calculated, incentive-driven decision to an internalized social habit once a critical tipping point is reached \citep{Ceschi2021, Corcoran2020}.

In this framework, the standard attitude update function, which typically suffers from a decay constant ($\delta$) due to enforcement fatigue or apathy, is refined by a dynamic damping factor ($D_{\mathrm{factor}}$). This modification aligns with recent findings on the persistence of pro-environmental behavior, suggesting that established social norms act as a buffer against the rapid erosion of compliance \citep{Andre2021}:

\begin{equation}
    A_{t+1} = A_t - (\delta \times D_{\mathrm{factor}})
\end{equation}

The damping factor is governed by the strength of local social norms ($SN$), representing the protective effect of community-wide compliance. Specifically, the model identifies a \textit{tipping point at 70\% compliance}, beyond which the behavior is considered internalized. This threshold reflects the ``critical mass'' theory in social dynamics, where minority behaviors cascade into majority conventions once a specific saturation point is exceeded \citep{Centola2018, Nyborg2024}:

\begin{equation}
    D_{\mathrm{factor}} = 
    \begin{cases} 
        0.1 & \text{if } SN > 0.70 \text{ (Strong Norms / Internalized)} \\
        0.5 & \text{if } SN > 0.50 \text{ (Moderate Norms)} \\
        1.0 & \text{otherwise (Weak Norms)}
    \end{cases}
\end{equation}

Furthermore, to prevent the total erosion of community standards during periods of fiscal austerity or low LGU intervention, the model introduces a \textit{Social Norm Floor set at 40\%}. This floor ensures that once a barangay has achieved a baseline level of awareness, the perceived social pressure does not drop to zero, representing the residual collective memory and existing communal infrastructure of the municipality \citep{Nishimura2022}. This refinement is critical for addressing the sustainability objectives of this research, as it allows the Deep Reinforcement Learning agent to discover strategies that foster self-sustaining compliance rather than perpetual reliance on high-cost monetary incentives .

\subsection{Household Decision-Making: The TPB Framework and Social Shields}


The behavioral state and daily segregation decisions of each Household Agent are governed by the Theory of Planned Behavior (TPB), a robust psychological framework that predicts deliberate behavior through the interplay of an individual's Attitude, Subjective Norms, and Perceived Behavioral Control (PBC) \citep{Meng2018, Zheng2022a}. However, traditional Agent-Based Models often suffer from hyper-reactivity, where agents immediately abandon positive behaviors the moment external pressures or local norms drop. 

To address this, we introduce the ``Two-Stage Social Shield'' mechanism, which mathematically models human behavioral inertia and the stickiness of pro-environmental habits \citep{Ceschi2021}. This mechanism dictates that once a community reaches critical thresholds of compliance, households develop a resistance to negative behavioral decay. \textit{Stage 1: Bayanihan Activation} is triggered when the barangay compliance surpasses 40\%, representing an emergent community spirit that slows the erosion of positive norms. \textit{Stage 2: Norm Lock-in} occurs when compliance exceeds 70\%, representing deep cultural habituation where segregation becomes the default state \citep{Jimenez2025}. In this locked-in state, households will largely maintain their segregation practices even if the perceived reality of neighbor compliance temporarily falters.

\begin{algorithm}[H]
\caption{Household Decision-Making (TPB with Shields)}
\label{alg:household_decision}
\begin{algorithmic}[1]
\renewcommand{\algorithmicrequire}{\textbf{Input:}}
\Require Household properties (Income, Base Attitude, Base SN, Base PBC)
\Statex
\State \Comment{Step 1: Update Social Norms (SN) with Perception Cap}
\State $Local\_Compliance \gets \text{Calculate\_Neighbors\_Compliance}(Radius = 2)$
\State $Authority\_Boost \gets Enforcement\_Intensity \times 0.50$
\State $Perceived\_Reality \gets \min(0.92, Local\_Compliance + Authority\_Boost)$
\Statex
\If{$Perceived\_Reality < Current\_SN$}
    \State \Comment{Apply Damper (Inertia/Stickiness)}
    \If{$Barangay\_Compliance > 0.70$}
        \State $Current\_SN \gets (0.90 \times Current\_SN) + (0.10 \times Perceived\_Reality)$ \Comment{Stage 2}
    \ElsIf{$Barangay\_Compliance > 0.40$}
        \State $Current\_SN \gets (0.85 \times Current\_SN) + (0.15 \times Perceived\_Reality)$ \Comment{Stage 1}
    \Else
        \State $Current\_SN \gets Perceived\_Reality$
    \EndIf
\Else
    \State $Current\_SN \gets Perceived\_Reality$
\EndIf
\State $Current\_SN \gets \min(0.95, Current\_SN)$ \Comment{Realism Cap}
\Statex
\State \Comment{Step 2: Update Attitude with Decay and Synergy}
\State $Decay\_Damper \gets \text{Calculate\_Damper\_Based\_On\_Compliance}()$
\State $Attitude \gets Attitude - (Decay\_Rate \times Decay\_Damper)$
\State $Attitude \gets Attitude + \text{Calculate\_IEC\_Synergy\_Boost}()$
\State $Attitude \gets \text{Apply\_Pride\_And\_Complacency\_Modifiers}(Attitude)$
\Statex
\State \Comment{Step 3: Calculate Utility and Decision}
\State $Monetary\_Impact \gets Incentive\_Value + (Fine\_Amount \times Prob\_Detection)$
\State $Net\_Cost \gets Base\_Effort - (Gamma\_Income \times (Monetary\_Impact / 2000.0))$
\State $Epsilon \gets \text{Random\_Gaussian}(\text{Mean} = 0, \text{StdDev} = 0.05)$
\Statex
\If{$Utility > 0.0$}
    \State $Is\_Compliant \gets \text{True}$
    \State \Comment{Human Error Rule}
    \If{$\text{Random}(0.0, 1.0) < 0.01$}
        \State $Is\_Compliant \gets \text{False}$
    \EndIf
\Else
    \State $Is\_Compliant \gets \text{False}$
\EndIf
\end{algorithmic}
\end{algorithm}

Following the parameter updates, the final household decision is computed using a randomized Utility maximization function. Mathematically, the $Utility$ equation integrates the weighted psychological components of the TPB with the objective economic costs, defined as:
\begin{equation}
    Utility = (w_a \cdot Attitude) + (w_{sn} \cdot Current\_SN) + (w_{pbc} \cdot Base\_PBC) - Net\_Cost + Epsilon
\end{equation}
where $w_a$, $w_{sn}$, and $w_{pbc}$ are the calibrated importance weights for each respective TPB construct, and $Epsilon$ represents the stochastic noise inherently present in human decision-making ($\mu=0, \sigma=0.05$). If the calculated $Utility$ exceeds $0.0$, the baseline intent is to segregate. 

However, intentionality does not guarantee perfect execution \citep{McCulloch2022}. To model this, we implement a strict ``1\% Human Error'' rule. Even when a household has a highly positive utility and full intention to comply, there is a 1\% probability that they fail to segregate on any given day due to unpredictable real-world factors (e.g., misplacing items, lack of time, or full bins). This constraint is functionally necessary to prevent the Agent-Based Model from plateauing at an unrealistic, deterministic 100\% compliance rate, ensuring the simulation maintains a degree of ecological validity aligned with actual municipal waste management data \citep{Tian2024}.

\subsection{Barangay and Enforcement Agent Design}

The \texttt{BarangayAgent} serves as the critical intermediate implementation layer, bridging the gap between municipal-level resource allocation and household-level behavioral shifts. The model initializes seven unique instances of this agent class—representing the geographically and socioeconomically distinct units of Liangan East, Esperanza, Poblacion, Binuni, Demologan, Mati, and Babalaya—calibrated with parameters derived from municipal institutional audits. Within this framework, the \texttt{BarangayAgent} functions as an executive entity that manages the ``stacking effect'' of fiscal resources, combining LGU-DRL quarterly augmentation funds with its own localized operating budget. 

Operating under the directive of the \texttt{BarangayAgent} is the \texttt{EnforcementAgent}, which serves as the operational extension of the municipal government tasked with physically monitoring household compliance and executing punitive measures.  Together, these agents mediate environmental policy by operationalizing the three core municipal levers: Information, Education, and Communication (IEC), Incentives, and Enforcement.

\subsubsection{IEC Campaign Broadcasting and Incentive Distribution}

When the LGU-DRL agent allocates IEC funds ($A_{\mathrm{IEC}, B_i}$) to a specific barangay, the \texttt{BarangayAgent} translates this continuous financial variable into a localized \textit{Information Dissemination Signal}. IEC campaigns are broadcast globally to all constituent \texttt{HouseholdAgent} objects within that barangay, proportionally shifting the internal Attitude weight ($w_A(t)$) positively to counteract the natural stochastic decay of pro-environmental habits.

For the Incentive budget ($A_{\mathrm{Inc}, B_i}$), the \texttt{BarangayAgent} functions as a direct financial distributor. During the quarterly evaluation phase, the agent identifies the subset of households actively practicing segregation (\texttt{is\_segregating == True}) and distributes the allocated reward pool among them, which directly modifies the perceived monetary benefit ($I$) in the household's net disutility calculation. Crucially, the \texttt{BarangayAgent} manages the per-capita dilution of these funds. If the local compliance rate surges, the fixed quarterly allocation is divided among a larger pool of households, reducing individual reward magnitude. This mechanism accurately simulates the ``victim of success'' dynamic, forcing the DRL agent to continuously adjust funding to maintain motivation.

\subsubsection{Enforcement Deployment, Capacity, and System Feedback}
To operationalize the ``No Segregation, No Collection'' mandate, the \texttt{BarangayAgent} deploys local \texttt{EnforcementAgent} staff. To ensure ecological validity, each enforcement agent is initialized with a \texttt{base\_daily\_capacity} of 30 households. This cap reflects the geographical realities and administrative friction of monitoring the municipality. Operating on a quarterly time step of 66 standard working days, the maximum theoretical inspection capacity ($C_{\mathrm{Quarterly}}$) of an active enforcer is defined as:

\begin{equation}
    C_{\mathrm{Quarterly}} = (C_{\mathrm{BaseDaily}} \times \mu_{\mathrm{eff}}) \times 66
    \label{eq:enforcement_capacity}
\end{equation}

\noindent where $\mu_{\mathrm{eff}}$ represents an efficiency multiplier (defaulted to 1.0) that can be adjusted to simulate varying levels of institutional competence.

During each simulation step, the deployed \texttt{EnforcementAgent} employs a randomized monitoring strategy, selecting a discrete sample of households up to its capacity limit. If a household is non-compliant, the agent flags it and increments its internal \texttt{citations\_issued} tracker, directly triggering the financial penalty ($F$) within the household's utility function. 

Crucially, the presence of these agents generates a systemic feedback loop. Following inspections, the enforcement agent updates the barangay's macro-level Perceived Probability of Detection ($P_{\mathrm{Detection}}$) using the ratio of total inspection capacity to the total household population:

\begin{equation}
    P_{\mathrm{Detection}} = \min\left(1.0, \frac{\sum C_{\mathrm{Quarterly}}}{N_{\mathrm{Households}}}\right)
    \label{eq:prob_detection}
\end{equation}

This updated probability is broadcasted back to the \texttt{HouseholdAgent} objects, dynamically increasing their anticipated risk. At the conclusion of each quarter, the \texttt{BarangayAgent} aggregates this data and reports the \textit{BarangayComplianceRate} to the LGU agent, completing a high-fidelity feedback loop where local citizen behavior informs subsequent municipal-level policy adjustments.

\subsection{Spatial Dynamics and Enforcement Constraints}


The physical environment of the Agent-Based Model (ABM) is constructed as a two-dimensional spatial grid, where Household and Enforcement Agents are geographically distributed to represent the topographical layout of the municipality. Unlike centralized, equation-based models that assume uniform mixing or omniscient actors, this spatial formulation enforces strict geographical limitations on agent interactions \citep{Meng2018}. Specifically, Enforcement Agents do not possess global, omniscient vision of the entire barangay. Instead, they operate under bounded rationality and physical constraints, forcing them to rely on localized spatial searches within a defined line-of-sight (e.g., $Radius = 10$). This architectural design is critical for the model's ecological validity, as it accurately simulates the limited monitoring capacity of real-world local enforcers who can only act upon the non-compliant behaviors they physically observe within their immediate vicinity \citep{Dalugdog2021, Tian2024}.

\begin{algorithm}[H]
\caption{Enforcement Agent Patrol \& Fining}
\label{alg:enforcement_patrol}
\begin{algorithmic}[1]
\renewcommand{\algorithmicrequire}{\textbf{Input:}}
\Require Agent Position, Local Radius = 10
\Statex
\State \Comment{Step 1: Mark Immediate Surroundings}
\State Mark all Household Agents in Radius 1 as ``Visited''
\Statex
\State \Comment{Step 2: Determine Movement (Separate Patrol Logic)}
\State $Local\_Targets \gets \text{Scan for Unvisited Household Agents in Radius 10}$
\State $\text{Filter } Local\_Targets \text{ for Non-Compliant behavior}$
\Statex
\If{$Local\_Targets \text{ is NOT empty}$}
    \State \Comment{CHASE MODE}
    \State $Target \gets \text{Find Closest Agent in } Local\_Targets$
    \State $Next\_Position \gets \text{Calculate coordinate step closer to } Target$
\Else
    \State \Comment{PATROL MODE}
    \State $Next\_Position \gets \text{Select Random Valid Adjacent Coordinate}$
\EndIf
\Statex
\State $\text{Move Agent to } Next\_Position$
\Statex
\State \Comment{Step 3: Enforcement Execution}
\State $Catch\_Zone \gets \text{Scan for Household Agents in Radius 1 (Include Center)}$
\For{each $Household$ in $Catch\_Zone$}
    \If{$Household \text{ is NOT Compliant}$}
        \State $\text{Apply Fine (Decrease Household Utility and Attitude)}$
        \State $\text{Increment Global Fines Collected}$
    \EndIf
\EndFor
\end{algorithmic}
\end{algorithm}


The algorithmic movement of the Enforcement Agents is bifurcated into two distinct behavioral states: ``Chase Mode'' and ``Patrol Mode.'' When in Patrol Mode, the agent conducts a stochastic spatial search, moving to random adjacent coordinates. This represents the routine, unguided neighborhood rounds conducted by personnel when no immediate waste violations are visible. However, if a non-compliant Household Agent is detected within the enforcer's scanning radius, the agent dynamically switches to Chase Mode, interrupting the random walk to calculate a directed spatial trajectory toward the closest violator to execute an intervention \citep{Zhao2022}. 

This spatial limitation profoundly reflects the real-world logistical constraints of municipal governance, particularly the deployment of \textit{barangay tanods} (village watchmen) or local environmental enforcers in the Philippines. Local Government Units (LGUs) rarely have the manpower or budgetary resources to monitor every household simultaneously \citep{Abordo2025}. Consequently, the probability of a household being caught and fined is not absolute; rather, it is a highly localized phenomenon dependent on spatial proximity and the enforcer's current patrol route. This localized enforcement mechanism organically generates the dynamic, fluctuating detection probability experienced by households, ensuring that the deterrent effect of punitive policies in the simulation is realistically constrained by the physical and temporal capacity of the barangay's law enforcement \citep{Jimenez2025}.